% ═══════════════════════════════════════════════════════════════════════════════
\chapter{The Causal Derivative: Diffusion and Resolution}
% ═══════════════════════════════════════════════════════════════════════════════
% CENTRAL CORRESPONDENCE:
%   Continuum: ∂f/∂x^μ, ∇_μ, □  (derivatives, connection, d'Alembertian)
%   Discrete: diffusion over the Hasse, local valence, net flux

\begin{introduction}
  \item The diffusion operator and the spectral dimension as a derivative.
  \item The Causal Resolution Theorem: precision of the discrete derivative.
  \item Flow gradients, inertia, and continuity equations.
\end{introduction}

If the Kuratowski integral gives us volume (the event count),
the derivative reveals dimensionality and curvature. In this
chapter, we define the differential operator not as a subtraction of
values at adjacent points, but as the rate of propagation of
information through the causal network.

% ─────────────────────────────────────────────────────────────────────────────
\section{The Diffusion Operator and the Spectral Dimension}

In Riemann calculus, the dimension of a manifold is given
a priori: one fixes $d = 4$ before writing a single
equation. In Kuratowski Calculus, the dimension is an
\emph{emergent}
property~\cite{carlip2017,ambjorn2005,eichhorn2014},
detectable via a diffusion operator (the random walker
on the Hasse diagram).

\begin{definition}[Random Walker on the Hasse Diagram]
\label{def:random-walker}
Let $H = (V, E)$ be the Hasse diagram (undirected) of a causal
set $\mathcal{C} = (V, \prec)$, where $E$ is the set of
covering links~$\covers$ with the orientation suppressed.
A \textbf{random walker} starts from an event~$u_0 \in V$ and, at
each discrete step $t \to t+1$, jumps to a neighbor chosen
uniformly at random:
\[
  \Pr\bigl[u_{t+1} = v \;\big|\; u_t = u\bigr]
  \;=\;
  \begin{cases}
    1/\deg(u) & \text{if } (u,v) \in E,\\[2pt]
    0 & \text{otherwise,}
  \end{cases}
\]
where $\deg(u) = |\{v : (u,v) \in E\}|$ is the total valence
(in-degree $+$ out-degree in the original DAG).
\end{definition}

\begin{definition}[Return Probability and Spectral Dimension]
\label{def:spectral-dimension}
The \textbf{return probability} $P(t)$ is the probability that
the walker returns to its origin after $t$~steps, averaged
over all possible origins:
\[
  P(t) \;=\; \frac{1}{|V|}\sum_{u \in V}
  \Pr\bigl[u_t = u \;\big|\; u_0 = u\bigr].
\]
The \textbf{spectral dimension} is extracted as the logarithmic
slope of this probability:
\begin{equation}
  \boxed{\;
  \dS(t) \;=\; -2\,\frac{d\log P(t)}{d\log t}.
  \;}
  \label{eq:spectral-dim}
\end{equation}
For a smooth $d$-dimensional manifold, $P(t) \sim t^{-d/2}$
and therefore $\dS(t) \to d$.
\end{definition}

\stoneref{$\dS$ is the \emph{inferential dimensionality}: the number
of independent inference axes of the deductive system.
$\dS \to 4$ means that there exist exactly four
orthogonal logical directions in which the system can reason. The
question ``why 4 dimensions?'' becomes ``why
4 independent inference axes?''}

\noindent
The spectral dimension is the fundamental \emph{derivative} of
Kuratowski Calculus: it measures how the propagation of
information decays with temporal depth. It is the discrete analogue
of the trace of the heat kernel
$K(t) = \mathrm{Tr}\,e^{-t\Delta}$ in Riemannian spectral
geometry~\cite{minakshisundaram1949,gilkey1975}. The
spectral dimension was introduced as a diagnostic for
causal sets in~\cite{eichhorn2014}; the phenomenon of
dimensional reduction at short scales ($\dS \to 2$) also appears
in causal dynamical
triangulations~\cite{ambjorn2005} and other approaches to quantum
gravity~\cite{carlip2009,modesto2009}.

\begin{theorem}[Dimensional Stabilization by Causal Prisms]
\phantomsection\label{thm:dimensional-stabilization}
The presence of Causal Prisms~$K_{2,N}$ in the Hasse diagram
\textbf{stabilizes} the spectral dimension $\dS$ around~$4$
at intermediate diffusion times.
\end{theorem}

\begin{proof}
A Causal Prism $K_{2,N}$ adds to the Hasse diagram a complete
bipartite subgraph: $N$ intermediaries, each connected
\emph{only} to both poles $u$ and $v$ (the intermediaries are
mutually disconnected by the triangle-free property). The
degree of the poles grows to $\deg(u), \deg(v) \geq N$, creating a
region of high connectivity around the poles.

The trapping mechanism operates via \textbf{polar ping-pong}:
a random walker (undirected) that reaches an intermediary
$w_i$ has exactly two neighbors ($u$ and $v$), so it jumps to
one of the poles with probability~$1/2$ each. Upon arriving at pole
$u$ (of degree $\geq N$), the walker chooses uniformly among its
$N$ intermediate children---with high probability it returns to the
belly of the prism instead of escaping into the vacuum. Thus, the
walker bounces between pole $\to$ intermediate $\to$ pole repeatedly,
generating a dwell time on the order of the \textbf{coupon
collector problem}: to visit all $N$ distinct intermediaries
requires $O(N\,H_N)$ steps, where $H_N =
\sum_{k=1}^{N} 1/k$.

This trapping modifies the return probability:
$P(t,u) = \Pr(X_t = u \mid X_0 = u)$ increases relative to the pure
vacuum at intermediate times, since the walker remains in the
neighborhood of the prism. Since $\dS = -2\,d(\ln P)/d(\ln t)$, an
increase in~$P$ in the range of diffusion times associated with the
prism scale reduces the logarithmic derivative, stabilizing $\dS$
near the value~$4$ corresponding to the dimension of the
macroscopic causal diamond ($d = 4$).

Without prisms, the pure vacuum (uniform Poisson) has degree
fluctuations of $O(\sqrt{\rho})$ that produce noisy variations
in~$P(t)$ and therefore in~$\dS$. The prisms smooth out these
fluctuations by imposing a regular connectivity structure.
\end{proof}

\begin{result}[title={Dimensional Stabilization by Matter}]
\phantomsection\label{res:dimensional-stabilization}
Matter is not a guest of four-dimensional spacetime; it is,
by operational definition, the topological structure that
\emph{forces} the derivative of information to behave
four-dimensionally
(Theorem~\ref{thm:dimensional-stabilization}).
\end{result}

\noindent\textit{Numerical observation.}\enspace
The simulation ($N$ up to $5 \times 10^6$) confirms that the pure
vacuum exhibits a noisy and fractal spectral dimension, while the
presence of prisms stabilizes $\dS \approx 4$ at intermediate
diffusion times, consistent with
Theorem~\ref{thm:dimensional-stabilization}.

% ─────────────────────────────────────────────────────────────────────────────
\section{The Causal Resolution Theorem}

The validity of any discrete derivative depends on the resolution
of the probe. In Kuratowski Calculus, the precision of our
description of spacetime is intrinsically tied to the
number of causal paths probed.

Theorem~\ref{thm:causal-resolution} (Chapter~6) establishes that
in order to maintain an error tolerance~$\epsilon$ in the measurement
of the geometry up to a depth~$t_{\max}$, the number of
walkers must satisfy:
\begin{equation}
  W \;\geq\; \frac{t_{\max}^{\,2}}{c\;\epsilon^2}.
  \tag{\ref{eq:causal-resolution}}
\end{equation}

\noindent
This result establishes the \textbf{certainty limit}:
the illusion of the Riemannian continuum is only sustainable if the
density of causal interactions is high enough to
average out the Poisson noise of the discrete topology. The scaling
$W \propto t_{\max}^2$ is a direct consequence of $\dS = 4$:
in a spacetime of spectral dimension~$d$, the general bound
is $W \propto t_{\max}^{d/2}$.

\begin{theorem}[General Spectral Resolution Bound]
\label{thm:resolution-general}
For a causal set with spectral dimension~$\dS$ at large
times, $P(t) \sim t^{-\dS/2}$. The resolution bound of
Theorem~\ref{thm:causal-resolution} generalizes to:
\begin{equation}
  W \;\geq\; \frac{t_{\max}^{\,\dS/2}}{c\;\epsilon^2}.
  \label{eq:resolution-general}
\end{equation}
In particular:
\begin{enumerate}
  \item $\dS = 2$ (UV limit, tree-like graphs):
        $W \propto t_{\max}$\quad (linear cost).
  \item $\dS = 4$ (macroscopic spacetime):
        $W \propto t_{\max}^2$\quad (quadratic cost).
  \item $\dS > 4$ (compact extra dimensions):
        the cost grows faster than quadratic.
\end{enumerate}
\end{theorem}

\begin{proof}
Follows identically the proof of
Theorem~\ref{thm:causal-resolution}, substituting
$P(t_{\max}) = c\,t_{\max}^{-\dS/2}$:
\[
  W \;\geq\; \frac{1}{\epsilon^2\,P(t_{\max})}
  \;=\; \frac{t_{\max}^{\,\dS/2}}{c\;\epsilon^2}.
\]
\end{proof}

\noindent
Formula~\eqref{eq:resolution-general} reveals that the
spectral dimension directly controls the
\emph{computational price of smoothness}: higher dimensions
are more costly to sustain.

% ─────────────────────────────────────────────────────────────────────────────
\section{Antichains: The Spatial Slices of the DAG}

To define spatial gradients over the causal network, we need
the notion of a ``discrete instant'': a cut of the DAG that separates
past from future.

\begin{definition}[Antichain]
\label{def:antichain}
An \textbf{antichain} is a subset $\mathcal{A} \subset V$
whose elements are mutually incomparable under the causal order:
\[
  \forall\, a, b \in \mathcal{A},\;a \neq b
  \;\colon\quad
  a \nprec b \;\;\text{and}\;\; b \nprec a.
\]
An antichain is the discrete analogue of a
\textbf{simultaneity hypersurface}: it defines an ``instant''
where no event can causally influence another.
\end{definition}

\noindent
The causal DAG admits a \textbf{Hasse stratum
foliation}~\cite{rideout2000,surya2019}: a partition
$V = \mathcal{A}_0 \cup \mathcal{A}_1 \cup \cdots \cup
\mathcal{A}_T$ where each stratum~$\mathcal{A}_k$ is a maximal
antichain, and every link $u \covers v$ connects a
stratum~$\mathcal{A}_k$ with the next stratum~$\mathcal{A}_{k+1}$.
The strata define a ``global discrete time'' $k = 0, 1,
\ldots, T$.

% ─────────────────────────────────────────────────────────────────────────────
\section{Flow Gradients and Inertia}

The Kuratowski derivative allows us to define force not as an
external vector field, but as a \emph{connectivity
gradient}.

\begin{definition}[Discrete d'Alembertian
Operator]
\label{def:box-operator}
Let $\phi \colon V \to \mathbb{R}$ be a scalar function over the
events of the causal set. The \textbf{discrete d'Alembertian}
(or causal wave operator) at an event~$u$ is defined as the
difference between the outgoing flux and the incoming flux through
the covering links:
\begin{equation}
  (\Box_{\mathcal{C}}\,\phi)(u)
  \;=\;
  \underbrace{\sum_{v :\, u \covers v} \phi(v)}_{\text{children of } u}
  \;-\;
  \underbrace{\sum_{w :\, w \covers u} \phi(w)}_{\text{parents of } u}.
  \label{eq:box-discrete}
\end{equation}
\end{definition}

\begin{figure}[ht]
\centering
\begin{tikzpicture}[
  event/.style={circle, fill=structure, inner sep=2pt, minimum size=6pt},
  layer1/.style={circle, fill=main, inner sep=1.5pt, minimum size=5pt},
  layer2/.style={circle, fill=third, inner sep=1.5pt, minimum size=5pt},
  layer3/.style={circle, fill=fourth, inner sep=1.5pt, minimum size=5pt},
  link/.style={draw=gray!40, thin},
  >=Stealth
]
% Central event
\node[event, label={[font=\small]above:$u$}] (u) at (0,0) {};
% Layer 1 (parents) -- coefficient +9
\foreach \i/\a in {1/200, 2/240, 3/280, 4/320} {
  \node[layer1] (L1\i) at (\a:1.5) {};
  \draw[link] (L1\i) -- (u);
}
\node[font=\scriptsize, main] at (260:2.1) {$L_1$: $+9$};
% Layer 2 (grandparents) -- coefficient -16
\foreach \i/\a in {1/195, 2/225, 3/255, 4/285, 5/315} {
  \node[layer2] (L2\i) at (\a:2.8) {};
  \draw[link, gray!25] (L2\i) -- (u);
}
\node[font=\scriptsize, third] at (255:3.4) {$L_2$: $-16$};
% Layer 3 (great-grandparents) -- coefficient +8
\foreach \i/\a in {1/210, 2/240, 3/270, 4/300} {
  \node[layer3] (L3\i) at (\a:4.0) {};
  \draw[link, gray!15] (L3\i) -- (u);
}
\node[font=\scriptsize, fourth] at (240:4.6) {$L_3$: $+8$};
% Central coefficient
\node[font=\scriptsize, structure] at (0.6,0.4) {$-1$};
% Temporal arrow
\draw[->, thick, structure!40] (2.5,0.5) -- (2.5,-4)
  node[midway, right, font=\small, color=structure!50] {past};
\end{tikzpicture}
\caption{The d'Alembertian template. The central event~$u$
(coefficient~$-1$) connects to three layers of causal ancestors. The
alternating-sign coefficients $\{-1,+9,-16,+8\}$ cancel the
volume of the interior and isolate the curvature signal at the
boundary---a discrete implementation of Huygens' principle.}
\label{fig:plantilla-dalembertiano}
\end{figure}

\noindent
This operator has a transparent interpretation. Let
$d^+(u) = |\{v : u \covers v\}|$ be the out-degree and
$d^-(u) = |\{w : w \covers u\}|$ the in-degree of~$u$. Then:
\begin{itemize}
  \item If $\phi \equiv 1$ (constant field):
    $(\Box_{\mathcal{C}}\,\phi)(u) = d^+(u) - d^-(u)$.
    The operator measures the \textbf{local causal asymmetry}---the
    difference between the number of immediate futures and immediate
    pasts.
  \item If $\phi$ varies, the operator detects how the
    signal \textbf{propagates}: the sum over children measures the
    dispersion toward the future, and the sum over parents measures the
    convergence from the past.
\end{itemize}

\noindent
Starting from~\eqref{eq:box-discrete}, the \textbf{normalized causal
gradient} in the temporal direction is:

\begin{definition}[Causal Gradient]
\label{def:causal-gradient}
\begin{equation}
  (\nabla_t\,\phi)(u)
  \;=\;
  \frac{1}{d^+(u)}\sum_{v :\, u \covers v} \phi(v)
  \;-\;
  \frac{1}{d^-(u)}\sum_{w :\, w \covers u} \phi(w).
  \label{eq:causal-gradient}
\end{equation}
It is the average toward the future minus the average from the past: the
discrete rate of change of~$\phi$ per topological step.
\end{definition}

\subsection{Inertia as Resistance to Flow}

When a walker traverses a Causal Prism~$K_{2,N}$, it experiences
a decrease in its effective propagation speed due to the
bifurcation of paths at the $N$ intermediate nodes. Let us formalize
this observation.

\begin{theorem}[Gradient Delay by Causal Prisms]
\label{thm:gradient-delay}
Let $\phi$ be a scalar field propagating through a
Prism~$\mathcal{P}(u,v,W)$ with $|W| = N$ intermediate nodes
(\S\ref{def:topological-resistance}). For the entry
pole~$u$, the causal gradient satisfies:
\begin{equation}
  (\nabla_t\,\phi)(u)
  \;=\;
  \frac{1}{N}\sum_{w \in W} \phi(w)
  \;-\;
  \frac{1}{d^-(u)}\sum_{w' :\, w' \covers u} \phi(w'),
  \label{eq:gradient-prism}
\end{equation}
where the first sum distributes uniformly over the $N$
intermediaries. If $\phi$ is constant over~$W$ (homogeneous field in
the belly of the prism), then:
\[
  (\nabla_t\,\phi)(u)
  \;=\; \phi_W - \bar{\phi}_{\text{pasado}},
\]
independently of~$N$. However, the information requires
$N$~diffusion steps to traverse the belly before reaching the
exit pole~$v$, producing an \textbf{effective delay}
$\Delta t_{\text{prisma}} = N$.
\end{theorem}

\begin{proof}
In the Prism $\mathcal{P}(u,v,W)$, the pole~$u$ has $d^+(u) = N$
children (all the $w \in W$). By
Definition~\ref{def:causal-gradient}:
\[
  (\nabla_t\,\phi)(u)
  \;=\;
  \frac{1}{N}\sum_{w \in W}\phi(w) - \bar{\phi}_{\text{pasado}}.
\]
A walker entering~$u$ must choose one of the $N$ intermediate
paths. Each intermediary~$w_i$ is connected to~$v$ by a
covering link $w_i \covers v$. But before the
information can propagate from~$u$ to~$v$, it must pass through at
least one intermediary: the minimum path is
$u \covers w_i \covers v$, of length~2. Compared with a
direct vacuum link $u \covers v'$ (length~1), the
Prism introduces a delay of $+1$ step for each intermediary
visited. If the walker explores the $N$ intermediaries by diffusion,
the mean first-passage time to pole~$v$ scales as $\Delta
t_{\text{prisma}} \sim N$.
\end{proof}

\noindent
This ``resistance'' to the flow of information is the mathematical
representation of inertia. The object does not ``weigh'' because it
possesses a mysterious intrinsic property; it weighs because the
derivative of its causal advance is delayed by the topology of the
Prism.

\begin{result}[title={The Derivative of Inertia}]
\phantomsection\label{res:inertia-derivative}
The topological mass $M = \kappa N$
(Axiom of Topological Inertia, Vol.~II) manifests
dynamically as the delay of the causal gradient:
\[
  \Delta t_{\text{prisma}} \;\sim\; N \;=\; M/\kappa.
\]
The greater the mass, the greater the delay. Inertia is, literally,
the number of covering links that information must process before
advancing one global step.
\end{result}

% ─────────────────────────────────────────────────────────────────────────────
\subsection{The Layered Operator (Benincasa--Dowker)}

The discrete d'Alembertian $\Box_{\mathcal{C}}$
(Def.~\ref{def:box-operator}) uses only the covering
links---the Hasse layer $L_1$.  To obtain second-order precision
in the continuum limit $\Box = \nabla_\mu \nabla^\mu$, we must
extend the sum to deeper layers of the Hasse diagram.

\begin{definition}[Hasse Layers]
Let $u$ be an event of the causal set.  The \textbf{Hasse layer}
$L_k(u)$ is the set of ancestors of~$u$ at Hasse distance
exactly~$k$:
\[
  L_k(u) \;=\; \bigl\{\,v : v \prec u
  \text{ at Hasse distance exactly } k\,\bigr\}.
\]
Thus, $L_1(u)$ are the direct parents, $L_2(u)$ the ``grandparents,'' etc.
\end{definition}

\begin{definition}[Layered Operator (Benincasa--Dowker)]
\phantomsection\label{def:box-layered}
The \textbf{layered operator} $\boxlayer$ is defined as
\begin{equation}
  \boxlayer\,\Phi(x)
  \;=\;
  \frac{4}{\lP^2\sqrt{6}}
  \Bigl(
    -\Phi(x)
    \;+\; 9\!\!\sum_{L_1}\!\Phi
    \;-\; 16\!\!\sum_{L_2}\!\Phi
    \;+\; 8\!\!\sum_{L_3}\!\Phi
  \Bigr).
  \label{eq:box-layered}
\end{equation}
\end{definition}

\noindent
The coefficients $\{-1, +9, -16, +8\}$ are not free parameters: they are
the unique solution of the system of three \textbf{moment conditions}
that require $\boxlayer$ to annihilate constant, linear,
and quadratic fields on a flat Poisson sprinkling.  The simple
operator $\Box_{\mathcal{C}}$ (Def.~\ref{def:box-operator}) is the
truncation to the layer $L_1$ alone.

\medskip\noindent\textbf{Sign alternation and complex phases.}\quad
The signs of the coefficients alternate: $(-,+,-,+)$.  This alternation
forces oscillatory behavior in the diffusion over the Hasse
graph.  In the continuum limit, preserving unitarity requires that
these oscillations fold into complex phases $e^{i\theta}$---this
is the discrete origin of the imaginary unit~$i$ in the
Schr\"odinger equation.  The complex Hilbert space of quantum
mechanics is not an independent axiom; it is a derived consequence
of the sign structure of~$\boxlayer$.

\begin{result}[title={The Layered Operator}]
The operator $\boxlayer$ has no free parameters: its coefficients
are dictated by the geometry.  The sign alternation forces
oscillatory diffusion that, in the continuum limit, produces the
complex phases of quantum mechanics.
\end{result}

% ─────────────────────────────────────────────────────────────────────────────
\subsection{Continuity Equations on the Graph}
\label{sec:discrete-continuity}

Flux conservation on a directed acyclic graph is not an
additional axiom: it is an algebraic consequence of the definition of
current over covering links~\cite{sorkin2012,johnston2009,aslanbeigi2014}.

\begin{definition}[Discrete Divergence]
\label{def:discrete-divergence}
Let $J : E_H \to \mathbb{R}$ be a \textbf{causal current} defined
over the edges of the Hasse diagram~$E_H$.  The \textbf{discrete
divergence} at an event~$u$ is
\begin{equation}
  (\operatorname{div}\,J)(u)
  \;=\;
  \sum_{v:\, u \covers v} J(u,v)
  \;-\;
  \sum_{w:\, w \covers u} J(w,u),
  \label{eq:discrete-divergence}
\end{equation}
where $\covers$ denotes the covering link (Hasse edge).
\end{definition}

\begin{theorem}[Discrete Continuity Equation]
\label{thm:discrete-continuity}
Let $J$ be a conserved current over~$E_H$.  For every event~$u$
in the interior of the causal set,
\[
  (\operatorname{div}\,J)(u) \;=\; 0.
\]
\end{theorem}

\begin{proof}
Consider an antichain~$\Sigma$ that separates the past from the
future of~$u$.  The total flux through~$\Sigma$ is
$\Phi(\Sigma) = \sum_{e \in E_H \cap \Sigma} J(e)$.
Since the Hasse diagram is a DAG, every edge entering~$u$
from~$\Sigma$ must exit through some edge toward the future of~$u$.
The telescoping sum over the antichain cancels term by
term: each incoming flux contribution appears exactly
once as outgoing flux. The result is
$\sum_{\text{in}} J - \sum_{\text{out}} J = 0$.
\end{proof}

\medskip\noindent\textbf{Continuum limit.}\quad
When $N \to \infty$ with density~$\rho = \lP^{-4}$, the convergence
of the Hasse diagram to the Lorentzian manifold
(Theorem~\ref{thm:emergent-metric}) carries the discrete divergence to the
classical operator: $\operatorname{div}\,J \to \partial_\mu J^\mu$.
The discrete continuity equation becomes
$\partial_\mu J^\mu = 0$, recovering the standard conservation law
of general relativity~\cite{aslanbeigi2014}.

\medskip\noindent\textbf{Connection with the random walker.}\quad
The random walker of the spectral dimension
(Def.~\ref{def:spectral-dimension}) is the probabilistic dual of the
deterministic current: its transition probabilities define a
stochastic current whose divergence vanishes by the conservation
of probability. Causal flux and spectral diffusion are two faces of
the same algebraic structure.

\begin{result}[title={Flux Conservation}]
Flux conservation on the DAG is not an additional axiom: it is an
algebraic consequence of the definition of current over covering
links. Every conserved current satisfies
$\operatorname{div}\,J = 0$ at each event, and this statement converges to
$\partial_\mu J^\mu = 0$ in the continuum limit.
\end{result}

\stoneref{Flux conservation is the consistency of the deductive
system: each premise produces exactly as many conclusions as
it receives from its own premises. A valid argument neither creates nor
destroys inferential information.}

% ███████████████████████████████████████████████████████████████████████████████
%   PART II : EMERGENT PHYSICS